% !TeX program = lualatex
% ================================================================
% Urdu Parallel Text Version of IntroducingAveragesStarters
%
% Generated by the server using the following settings:
%
%   language            Urdu (urd)
%   text direction      right to left
%   font                Noto Nastaliq Urdu
%   fallback font       Noto Serif
%   built from          latex/starters/targeted/KS3/statistics/introducingAveragesStarters.tex
%   vocabulary key      latex/starters/targeted/KS3/statistics/introducingAveragesStarters/L2/urd/introducingAveragesStarters_urduKey.tex
%   tier 3 vocabulary   green   RGB 0 128 0
%   English text        black   RGB 0 0 0
%   Urdu text           purple  RGB 112 48 160
%   layout macros       version 3
%
% Please don't edit any information in this comment block - the server
% compares it against the current settings and rebuilds this file when
% they differ, so changes here would be overwritten. However, feel free
% to make updates to the LaTeX below if you want to edit it.
% ================================================================

\documentclass{beamer}
% ================================================================
% Copied from the English sheet this was translated from, so that
% anything it sets up for itself still works here
% ================================================================
\usepackage{tikz}
\usetikzlibrary{calc,arrows.meta,patterns}
\usepackage{multicol}
\usepackage{amsmath}

% ================================================================
% Answer-overlay helpers
% ================================================================
\newcommand{\ablank}[1]{%
  \alt<2>{\textcolor{red}{#1}}{\underline{\phantom{#1}}}%
}
\newcommand{\ashow}[1]{\uncover<2->{\textcolor{red}{\small #1}}}
\newcommand{\ashowq}[1]{\alt<2->{\textcolor{red}{\small #1}}{?}}

% Answers drawn straight on to a TikZ diagram, revealed on overlay 2 like \ashow.
% Space is reserved on overlay 1, so the diagram does not move between slides.
\usetikzlibrary{overlay-beamer-styles}
\tikzset{
  ans/.style     = {red, thick, visible on=<2->},
  ansfill/.style = {red!15, visible on=<2->},
  anslab/.style  = {red, font=\tiny, visible on=<2->},
}

\title{Averages: Mean \& Median}
\author{}
\date{}

% Lego tower: \legotower{x-offset}{height}{colour}
% Each brick is 0.65 wide x 0.55 tall, with two studs
\newcommand{\legotower}[3]{%
  \foreach \i in {1,...,#2}{%
    \draw[fill=#3!55, draw=#3!80!black, line width=0.35pt]
      (#1, {(\i-1)*0.55}) rectangle (#1+0.65, {\i*0.55});
    \draw[fill=#3!80, draw=#3!80!black, line width=0.2pt]
      (#1+0.15,{(\i-1)*0.55+0.38}) circle (0.08);
    \draw[fill=#3!80, draw=#3!80!black, line width=0.2pt]
      (#1+0.44,{(\i-1)*0.55+0.38}) circle (0.08);
  }%
}

  \newcommand{\lego}[3]{%
    \foreach \ii in {1,...,#2}{%
      \draw[fill=#3!55,draw=#3!80!black,line width=0.3pt]
        (#1,{(\ii-1)*0.42}) rectangle (#1+0.6,{\ii*0.42});
      \draw[fill=#3!80,draw=#3!80!black,line width=0.15pt]
        (#1+0.13,{(\ii-1)*0.42+0.30}) circle (0.07);
      \draw[fill=#3!80,draw=#3!80!black,line width=0.15pt]
        (#1+0.40,{(\ii-1)*0.42+0.30}) circle (0.07);
    }%
  }


  \newcommand{\legob}[3]{%
    \foreach \ii in {1,...,#2}{%
      \draw[fill=#3!55,draw=#3!80!black,line width=0.3pt]
        (#1,{(\ii-1)*0.42}) rectangle (#1+0.6,{\ii*0.42});
      \draw[fill=#3!80,draw=#3!80!black,line width=0.15pt]
        (#1+0.13,{(\ii-1)*0.42+0.30}) circle (0.07);
      \draw[fill=#3!80,draw=#3!80!black,line width=0.15pt]
        (#1+0.40,{(\ii-1)*0.42+0.30}) circle (0.07);
    }%
  }

% Characters this language's font has not got are borrowed from another,
% rather than being silently dropped - see docs/translated-sheet-latex.md
\directlua{%
  if luaotfload and luaotfload.add_fallback then
    luaotfload.add_fallback("eallatinfallback",
      { "Noto Serif:mode=harf;" })
  else
    tex.error("this luaotfload is too old to register a fallback font")
  end
}
\usepackage[english,bidi=basic]{babel}
\babelprovide[import]{urdu}
\babelfont[urdu]{rm}[Renderer=Harfbuzz, BoldFont={Noto Nastaliq Urdu}, ItalicFont={Noto Nastaliq Urdu}, BoldItalicFont={Noto Nastaliq Urdu}, RawFeature={fallback=eallatinfallback}]{Noto Nastaliq Urdu}
\babelfont[urdu]{sf}[Renderer=Harfbuzz, BoldFont={Noto Nastaliq Urdu}, ItalicFont={Noto Nastaliq Urdu}, BoldItalicFont={Noto Nastaliq Urdu}, RawFeature={fallback=eallatinfallback}]{Noto Nastaliq Urdu}
\newcommand{\ealtext}[1]{\foreignlanguage{urdu}{#1}}
\newcommand{\ealtextblock}[1]{{\raggedleft\foreignlanguage{urdu}{#1}\par}}
% ================================================================
% EAL layout helpers
% ================================================================
\usepackage{xcolor}

\definecolor{ealkeycolour}{RGB}{0,128,0}
\definecolor{ealenglish}{RGB}{0,0,0}
\definecolor{ealtwocolour}{RGB}{112,48,160}

\newsavebox{\ealboxen}
\newsavebox{\ealboxtr}
\newlength{\ealglosswd}
\newlength{\ealglossraise}
\setlength{\ealglossraise}{1.15em}

% a tier 3 word, in English and in translation alike
\newcommand{\ealkey}[1]{\textcolor{ealkeycolour}{#1}}
\newcommand{\ealkeytr}[1]{\textcolor{ealkeycolour}{\ealtext{#1}}}

% One translated word above one English word. Built from kernel
% primitives, never a tabular, which would break wherever the sheet
% loads array, booktabs or tabularx. The box is as wide as the wider of
% the two, so a long translation pushes its neighbours aside instead of
% overlapping them, and it claims its full height so a table row or a
% TikZ node makes room for it.
\newcommand{\ealgloss}[2]{%
  \sbox{\ealboxen}{\ealkey{#1}}%
  \sbox{\ealboxtr}{\small\textcolor{ealtwocolour}{\ealtext{#2}}}%
  \setlength{\ealglosswd}{\wd\ealboxen}%
  \ifdim\wd\ealboxtr>\ealglosswd\setlength{\ealglosswd}{\wd\ealboxtr}\fi
  \leavevmode
  \makebox[\ealglosswd][c]{%
    \makebox[0pt][c]{\usebox{\ealboxen}}%
    \makebox[0pt][c]{\raisebox{\ealglossraise}{\usebox{\ealboxtr}}}%
  }%
}

% opens up line spacing around a block of glosses, so the raised
% translations sit clear of the line above
\newenvironment{ealglossed}{\par\linespread{2.1}\selectfont\ignorespaces}{\par}

% a whole translated sentence above its English counterpart. block level,
% so it does not belong inside a tabular cell
\newcommand{\ealpara}[2]{%
  \par\addvspace{0.5\baselineskip}%
  {\color{ealtwocolour}\ealtextblock{#1}}%
  \penalty200
  {\color{ealenglish}#2\par}%
  \addvspace{0.5\baselineskip}%
}

\begin{document}

\begin{frame}
\ealpara{\ealkeytr{اوسط}: \ealkeytr{اوسط} اور \ealkeytr{میڈین}}{\ealkey{Averages}: \ealkey{Mean} \& \ealkey{Median}}
\end{frame}

%------------------------------------------------
% STARTER 1
%------------------------------------------------
\begin{frame}{Starter}
\ealpara{اسٹارٹر}{Starter}
\vspace{-0.8em}
\begin{columns}[T]

\column{0.33\textwidth}
\begin{center}

\ealpara{ان کو متوازن کریں}{Q7: balance these}
\begin{tikzpicture}[scale=0.9]
  \legotower{0.00}{1}{blue}
  \legotower{0.75}{3}{blue}
  \legotower{1.50}{2}{blue}
  \node[below,font=\tiny] at (0.325,0) {1};
  \node[below,font=\tiny] at (1.075,0) {3};
  \node[below,font=\tiny] at (1.825,0) {2};
  \node[below,font=\tiny\bfseries] at (1.0,{-0.38})
    {\textbf{Q7}: balance these};
  \draw[ans, dashed] (-0.15,1.1) -- (2.15,1.1);
  \node[anslab, right] at (2.15,1.1) {2};
\end{tikzpicture}

\vspace{0.55em}

\ealpara{ان کو متوازن کریں}{Q8: balance these}
\begin{tikzpicture}[scale=0.78]
  \legotower{0.00}{1}{orange}
  \legotower{0.75}{3}{orange}
  \legotower{1.50}{5}{orange}
  \legotower{2.25}{3}{orange}
  \legotower{3.00}{3}{orange}
  \node[below,font=\tiny] at (0.325,0) {1};
  \node[below,font=\tiny] at (1.075,0) {3};
  \node[below,font=\tiny] at (1.825,0) {5};
  \node[below,font=\tiny] at (2.575,0) {3};
  \node[below,font=\tiny] at (3.325,0) {3};
  \node[below,font=\tiny\bfseries] at (1.7,{-0.38})
    {\textbf{Q8}: balance these};
  \draw[ans, dashed] (-0.15,1.65) -- (3.75,1.65);
  \node[anslab, right] at (3.75,1.65) {3};
\end{tikzpicture}

\end{center}

\column{0.67\textwidth}
\footnotesize
\vspace{0.2em}
\begin{enumerate}\itemsep1pt
  \begin{multicols}{2}
    \item \ealpara{12 ÷ 4 = 3}{$12 \div 4 = \ashowq{3}$}
    \item \ealpara{18 ÷ 3 = 6}{$18 \div 3 = \ashowq{6}$}
  \end{multicols}
  \begin{multicols}{2}
    \item \ealpara{15 ÷ 3 = 5}{$15 \div 3 = \ashowq{5}$}
    \item \ealpara{28 ÷ 4 = 7}{$28 \div 4 = \ashowq{7}$}
  \end{multicols}
  \item \ealpara{چار دوست 20 مٹھائیاں برابر تقسیم کرتے ہیں۔ ہر ایک کو کتنی ملیں گی؟}{Four friends share 20 sweets equally. How many do they each get? \ashow{5 sweets}}
  \item \ealpara{تین بچوں کے پاس 2، 5 اور 8 اسٹیکرز ہیں۔ وہ انہیں ملا کر برابر تقسیم کرتے ہیں۔ ہر بچے کو کتنے ملیں گے؟}{Three children have 2, 5 and 8 stickers. They combine them and share them equally. How many does each child get? \ashow{5 stickers}}
  \item \ealpara{اوپر والے ٹاورز (بائیں) کی تصویر بنائیں، جو اینٹوں کی کل تعداد تبدیل کیے بغیر ایک ہی اونچائی تک متوازن ہوں۔}{Draw a picture of the top towers (left), balanced to the same height without changing the total number of bricks.}
  \item \ealpara{نیچے والے ٹاورز (بائیں) کی تصویر بنائیں، جو اینٹوں کی کل تعداد تبدیل کیے بغیر ایک ہی اونچائی تک متوازن ہوں۔}{Draw a picture of the bottom towers (left), balanced to the same height without changing the total number of bricks.}
\end{enumerate}

\end{columns}
\end{frame}

%------------------------------------------------
% STARTER 2
%------------------------------------------------
\begin{frame}{Starter}
\ealpara{اسٹارٹر}{Starter}
\vspace{-0.8em}
\begin{columns}[T]

\column{0.33\textwidth}
\begin{center}

\ealpara{ان کو متوازن کریں}{Q7: balance these}
\begin{tikzpicture}[scale=0.85]

  \lego{0.00}{3}{blue}
  \lego{0.70}{5}{blue}
  \lego{1.40}{4}{blue}
  \lego{2.10}{8}{blue}
  \node[below,font=\tiny] at (0.30,0) {3};
  \node[below,font=\tiny] at (1.00,0) {5};
  \node[below,font=\tiny] at (1.70,0) {4};
  \node[below,font=\tiny] at (2.40,0) {8};
  \node[below,font=\tiny\bfseries] at (1.35,{-0.38})
    {\textbf{Q7}: balance these};
  \draw[ans, dashed] (-0.15,2.10) -- (2.85,2.10);
  \node[anslab, right] at (2.85,2.10) {5};
\end{tikzpicture}

\vspace{0.45em}

\ealpara{ان کو متوازن کریں؟}{Q8: balance these?}
\begin{tikzpicture}[scale=0.78]

  \legob{0.00}{1}{orange}
  \legob{0.75}{5}{orange}
  \legob{1.50}{2}{orange}
  \legob{2.25}{8}{orange}
  \legob{3.00}{4}{orange}
  \node[below,font=\tiny] at (0.30,0) {1};
  \node[below,font=\tiny] at (1.05,0) {5};
  \node[below,font=\tiny] at (1.80,0) {2};
  \node[below,font=\tiny] at (2.55,0) {8};
  \node[below,font=\tiny] at (3.30,0) {4};
  \node[below,font=\tiny\bfseries] at (1.7,{-0.38})
    {\textbf{Q8}: balance these?};
  \draw[ans, dashed] (-0.15,1.68) -- (3.75,1.68);
  \node[anslab, right] at (3.75,1.68) {4};
\end{tikzpicture}

\end{center}

\column{0.67\textwidth}
\footnotesize
\vspace{0.2em}
  \begin{multicols}{2}
  \begin{enumerate}
    \item \ealpara{7 × 8 = 56}{$7 \times 8 = \ashowq{56}$}
    \item \ealpara{63 ÷ 7 = 9}{$63 \div 7 = \ashowq{9}$}
    \end{enumerate}
  \end{multicols}
  \begin{multicols}{2}
  \begin{enumerate}
  \setcounter{enumi}{2}
    \item \ealpara{48 ÷ 8 = 6}{$48 \div 8 = \ashowq{6}$}
    \item \ealpara{5 × 9 = 45}{$5 \times 9 = \ashowq{45}$}
    \end{enumerate}
  \end{multicols}
  \begin{enumerate}
  \setcounter{enumi}{4}
  \item \ealpara{پانچ بچوں کے پاس 4، 7، 3، 6، 5 ماربلز ہیں۔ برابر تقسیم کریں تو ہر ایک کو کتنے ملیں گے؟}{Five children have 4, 7, 3, 6, 5 marbles. Shared equally, how many each? \ashow{5 marbles}}
  \item \ealpara{چھ دوست کار واش پر £8، £14، £10، £6، £12، £10 کماتے ہیں۔ منصفانہ تقسیم کریں — ہر ایک کو کتنا ملے گا؟}{Six friends earn £8, £14, £10, £6, £12, £10 at a car wash. Split fairly — how much each? \ashow{\pounds10 each}}
  \item \ealpara{آٹھ رنرز 54، 48، 61، 57، 50، 63، 55، 52 سیکنڈ میں دوڑ مکمل کرتے ہیں۔ بغیر کوئی حساب کیے، \ealkeytr{اوسط} وقت کا \ealkeytr{تخمینہ} لگائیں۔}{Eight runners finish in 54, 48, 61, 57, 50, 63, 55, 52 s. Without doing any calculations, \ealkey{estimate} the \ealkey{average} time. \ashow{about 55 s}}
  \item \ealpara{اوپر والے ٹاورز کس اونچائی تک متوازن ہوں گے؟}{What heights would the top towers balance out to? \ashow{5 bricks high}}
  \item \ealpara{نیچے والے ٹاورز کس اونچائی تک متوازن ہوں گے؟}{What heights would the bottom towers balance out to? \ashow{4 bricks high}}
\end{enumerate}

\end{columns}
\end{frame}

%------------------------------------------------
% STARTER 3
%------------------------------------------------
\begin{frame}{Starter}
\ealpara{اسٹارٹر}{Starter}
\vspace{-0.8em}

\footnotesize
\vspace{0.2em}
\begin{enumerate}\itemsep1pt
  \begin{multicols}{2}
    \item \ealpara{3 + 4 + 7 = 14}{$3 + 4 + 7 = \ashowq{14}$}
    \item \ealpara{15+27+18=60}{$15+27+18=\ashowq{60}$}
    \item \ealpara{60÷4=15}{$60\div 4=\ashowq{15}$}
    \item \ealpara{-15 ÷ 3=-5}{$-15 \div 3=\ashowq{-5}$}
  \end{multicols}
  \vspace{-1em}
  \item \ealpara{$4، 6، 2، 8$ کا \ealkeytr{اوسط} کیا ہے؟}{What is the \ealkey{mean} of: $4,\ 6,\ 2,\ 8$ \ashow{5}}
  \item \ealpara{$3، 9، 5، 7، 1$ کا \ealkeytr{اوسط} کیا ہے؟}{What is the \ealkey{mean}  of: $3,\ 9,\ 5,\ 7,\ 1$ \ashow{5}}
  \item \ealpara{$-3، 5، -1، 7، 2$ کا \ealkeytr{اوسط} کیا ہے؟}{What is the \ealkey{mean}  of: $-3,\ 5,\ -1,\ 7,\ 2$ \ashow{2}}
  \item \ealpara{$-8، -2، -5، -1، -4$ کا \ealkeytr{اوسط} کیا ہے؟}{What is the \ealkey{mean}  of: $-8,\ -2,\ -5,\ -1,\ -4$ \ashow{$-4$}}
\end{enumerate}

\begin{enumerate}\itemsep1pt
  \setcounter{enumi}8
  \item \ealpara{چار اعداد کا \ealkeytr{اوسط} 9 ہے۔ تین اعداد 6، 11 اور 8 ہیں۔ چوتھا عدد تلاش کریں۔}{The \ealkey{mean} of four numbers is 9. Three are 6, 11 and 8. Find the fourth. \ashow{11}}
  \item \ealpara{$3، 7، x، 10$ کا \ealkeytr{اوسط} $8$ ہے۔ $x$ کیا ہے؟}{The \ealkey{mean} of $3,\ 7,\ x,\ 10$ is $8$. What is $x$? \ashow{12}}
\end{enumerate}


\end{frame}

%------------------------------------------------
% STARTER 4
%------------------------------------------------
\begin{frame}{Starter 4}
\ealpara{اسٹارٹر 4}{Starter 4}
\vspace{-0.8em}

\begin{columns}[T]

% ---------------- LEFT COLUMN ----------------
\column{0.33\textwidth}

\begin{center}

\vspace{7em}

\ealpara{Q10: سڑک پر گھروں کی قیمتیں}{Q10: house prices on a road}
\begin{tikzpicture}[scale=0.72]

  \draw[thick,->] (-0.2,0) -- (4.6,0);

  \foreach \xd in {0.3,0.55,0.8,1.05,1.3}{
    \fill[blue] (\xd,0.22) circle (0.1);
  }

  \fill[red] (4.2,0.22) circle (0.1);

  \node[below,font=\tiny,blue] at (0.8,-0.1)
    {\pounds125--140k};

  \node[below,font=\tiny,red] at (4.2,-0.1)
    {\pounds980k};

  \draw[<->,dashed,gray]
    (1.35,0.52) -- (4.1,0.52)
    node[midway,above,font=\tiny,gray]{big gap};

  \node[below,font=\tiny\bfseries]
    at (2.1,-0.55)
    {\textbf{Q10}: house prices on a road};

\end{tikzpicture}

\end{center}

% ---------------- RIGHT COLUMN ----------------
\column{0.67\textwidth}

\footnotesize
\vspace{0.2em}

% First four quick questions
\begin{multicols}{2}
\begin{enumerate}\itemsep1pt
  \item \ealpara{4+2+10=16}{$4+2+10=\ashowq{16}$}
  \item \ealpara{-4+9=5}{$-4+9=\ashowq{5}$}
  \item \ealpara{-3-5=-8}{$-3-5=\ashowq{-8}$}
  \item \ealpara{6×4=24}{$6\times4=\ashowq{24}$}
\end{enumerate}
\end{multicols}

% Remaining questions
\begin{enumerate}\itemsep1pt
\setcounter{enumi}{4}

  \item \ealpara{چھوٹے سے بڑے تک \ealkeytr{ترتیب دیں}: $8,3,11,5,7$۔ درمیانی \ealkeytr{قدر} کون سی ہے؟}{Order smallest to largest: $8,3,11,5,7$. Which is the middle \ealkey{value}? \ashow{7}}

  \item \ealpara{\ealkeytr{ترتیب دیں}: $21,14,9,35,17,28,6$۔ درمیانی \ealkeytr{قدر} کون سی ہے؟}{Order: $21,14,9,35,17,28,6$. Which is the middle \ealkey{value}? \ashow{17}}

  \item \ealpara{$3,7,8,9,41$ کے لیے: \ealkeytr{تخمینہ} لگائیں کہ درمیانی \ealkeytr{قدر} \ealkeytr{اوسط} سے اوپر ہے یا نیچے۔ پھر جانچیں۔}{For $3,7,8,9,41$: \textit{\ealkey{estimate}} whether the middle \ealkey{value} is above or below the \ealkey{mean}. Then check. \ashow{below: \ealkey{median} $8$, \ealkey{mean} $13.6$}}

  \item \ealpara{درمیانی \ealkeytr{قدر} کیا ہے: $-7,\ -2,\ -5,\ 1,\ -1,\ 3,\ -3$}{What is the middle \ealkey{value} of: $-7,\ -2,\ -5,\ 1,\ -1,\ 3,\ -3$ \ashow{$-2$}}

  \item \ealpara{درمیانی \ealkeytr{قدر} کیا ہے: $\tfrac{1}{2},\ \tfrac{1}{4},\ \tfrac{3}{4},\ \tfrac{1}{8},\ \tfrac{5}{8}$}{What is the middle \ealkey{value} of: $\tfrac{1}{2},\ \tfrac{1}{4},\ \tfrac{3}{4},\ \tfrac{1}{8},\ \tfrac{5}{8}$ \ashow{$\frac{1}{2}$}}

  \item \ealpara{ایک سڑک پر گھروں کی قیمتیں یہ ہیں: \pounds130k, \pounds125k, \pounds140k, \pounds128k, \pounds135k, \textcolor{red}{\pounds980k}۔ \ealkeytr{اوسط} کیوں کارآمد نہیں ہو سکتا؟}{House prices on a road are: \pounds130k, \pounds125k, \pounds140k, \pounds128k, \pounds135k, \textcolor{red}{\pounds980k}. Why might the \ealkey{mean} not be helpful? \ashow{the \pounds980k \ealkey{outlier} pulls the \ealkey{mean} up}}

\end{enumerate}

\end{columns}

\end{frame}

%------------------------------------------------
% STARTER 5
%------------------------------------------------
\begin{frame}{Starter}
\ealpara{اسٹارٹر}{Starter}
\vspace{-0.8em}
\begin{columns}[T]

\column{0.33\textwidth}
\begin{center}

\vspace{7em}

\ealpara{\ealkeytr{اوسط} \ealkeytr{آؤٹ لیرز} کی وجہ سے کھنچ جاتا ہے}{\ealkey{Mean} is pulled by \ealkey{outliers}}
\begin{tikzpicture}[scale=0.72]
  \draw[thick,->] (-0.2,0) -- (4.6,0);
  \foreach \xd in {0.8, 1.05, 1.3, 1.55}{
    \fill[blue] (\xd,0.22) circle (0.1);
  }
  \fill[red] (4.1,0.22) circle (0.1);
  \node[below,font=\tiny,blue] at (1.2,{-0.1})  {£24--28k};
  \node[below,font=\tiny,red]  at (4.1,{-0.1})  {£97k};
  \draw[blue,very thick] (1.175,0.42) -- (1.175,-0.08);
  \node[above,blue,font=\tiny] at (1.175,0.44) {median};
  \draw[->,thick,orange] (1.9,0.22) -- (2.55,0.22);
  \node[above,orange,font=\tiny] at (2.25,0.26) {mean};
  \node[below,font=\tiny\bfseries] at (2.0,{-0.55})
    {Mean is pulled by outliers};
\end{tikzpicture}

\vspace{0.5em}


\end{center}

\column{0.67\textwidth}
\footnotesize
\vspace{0.3em}
\begin{enumerate}\itemsep1pt
  \begin{multicols}{2}
    \item \ealpara{-6+(-2)=-8}{$-6+(-2) = \ashowq{-8}$}
    \item \ealpara{1/3+1/2=5/6}{$\frac{1}{3}+\frac{1}{2} = \ashowq{\frac{5}{6}}$}
  \end{multicols}
  \vspace{-1em}
  \item \ealpara{$3، 9، 4، 7، 1$ کا \ealkeytr{میڈین}}{Median of: $3,\ 9,\ 4,\ 7,\ 1$ \ashow{4}}
  \item \ealpara{$5، 8، 3، 12، 7$ کا \ealkeytr{اوسط} اور \ealkeytr{میڈین}۔ کون سا بڑا ہے؟}{\ealkey{Mean} \textbf{and} \ealkey{median} of: $5,\ 8,\ 3,\ 12,\ 7$. Which is larger? \ashow{equal (both $7$)}}
  \item \ealpara{پانچ تنخواہیں یہ ہیں: £24k, £26k, £25k, £28k, \textcolor{red}{£97k}۔ کیا آپ \ealkeytr{اوسط} معلوم کرنے کے لیے \ealkeytr{اوسط} یا \ealkeytr{میڈین} استعمال کرنا پسند کریں گے؟}{Five salaries are: £24k, £26k, £25k, £28k, \textcolor{red}{£97k}. Would you prefer to use the \ealkey{mean} or the \ealkey{median} to find the \ealkey{average}? \ashow{the median}}
  \item \ealpara{$-10,-1,0,2,3$ کے لیے: \ealkeytr{تخمینہ} لگائیں کہ \ealkeytr{اوسط} \ealkeytr{مثبت} ہے یا \ealkeytr{منفی}۔ پھر \ealkeytr{اوسط} اور \ealkeytr{میڈین} دونوں معلوم کریں۔}{For $-10,-1,0,2,3$: \textit{\ealkey{estimate}} whether the \ealkey{mean} is \ealkey{positive} or \ealkey{negative}. Then calculate both \ealkey{mean} and \ealkey{median}. \ashow{negative: \ealkey{mean} $-1.2$; \ealkey{median} $0$}}
  \item \ealpara{$-6، -2، 4، -3، 7، 0$ کا \ealkeytr{اوسط} اور \ealkeytr{میڈین} معلوم کریں۔}{Find the \ealkey{mean} and \ealkey{median} of: $-6,\ -2,\ 4,\ -3,\ 7,\ 0$ \ashow{\ealkey{mean} $0$; \ealkey{median} $-1$}}
  \item \ealpara{$\frac{1}{3},\ \frac{3}{4},\ \frac{1}{2},\ \frac{1}{6},\ \frac{2}{3},\ \frac{1}{4}$ کا \ealkeytr{میڈین} کیا ہے؟}{What is the \ealkey{median} of: $\frac{1}{3},\ \frac{3}{4},\ \frac{1}{2},\ \frac{1}{6},\ \frac{2}{3},\ \frac{1}{4}$ \ashow{$\frac{5}{12}$}}
  \item \ealpara{صحیح یا غلط؟\\ (a) \ealkeytr{میڈین} ہمیشہ \ealkeytr{ڈیٹا سیٹ} کی \ealkeytr{اقدار} میں سے ایک ہوتا ہے۔\\ (b) \ealkeytr{اوسط} اور \ealkeytr{میڈین} ہمیشہ مختلف ہوتے ہیں۔\\ (c) بہت بڑا \ealkeytr{آؤٹ لیر} \ealkeytr{میڈین} کو \ealkeytr{اوسط} سے زیادہ متاثر کرتا ہے۔}{True or False?\\
    (a) The \ealkey{median} is always one of the \ealkey{values} in the \ealkey{dataset}. \ashow{False}\\
    (b) The \ealkey{mean} and \ealkey{median} are always different. \ashow{False}\\
    (c) A very large \ealkey{outlier} affects the \ealkey{median} more than the \ealkey{mean}. \ashow{False}}

\end{enumerate}

\end{columns}
\end{frame}

\end{document}